%برای شروع به فایل help.pdf رجوع کنید و فقط توجه داشته باشید کلاس VThesis.cls به کلاس SBU.cls، فایل اصلی از VThesis.tex به main.tex و پوشه از VThesis_Tabriz به AMThesis_Beheshti تغییر نام داده‌اند.
\documentclass[a4paper,11pt,msc]{SBU}
% در این فایل، دستورها و تنظیمات مورد نیاز، آورده شده است.
%-------------------------------------------------------------------------------------------------------------------

\usepackage{stmaryrd}
\usepackage{amssymb,amsmath,amsthm}
%ایجاد فرورفتگی در اولین پاراگراف هر بخش
\usepackage{indentfirst}
% بسته‌ای برای تنطیم حاشیه‌های بالا، پایین، چپ و راست صفحه
\usepackage[top=30mm, bottom=30mm, left=25mm, right=35mm]{geometry}
% بسته‌‌ای برای ظاهر شدن شکل‌ها و تصاویر متن
\usepackage{graphicx}
% بسته‌ای برای رسم کادر
\usepackage{framed} 
% بسته‌‌ای برای چاپ شدن خودکار تعداد صفحات در صفحه «معرفی پایان‌نامه»
\usepackage{lastpage}
% بسته‌‌ای برای ایجاد دیاگرام‌های مختلف
\usepackage[all]{xy}
\usepackage{tikz}
\usetikzlibrary{topaths}
\tikzset{terminal/.style={
							% The shape:
							rectangle ,minimum size =6mm,rounded corners=4mm,
							% The rest
							%very thick,draw=black!50,
							top color=green!10,bottom color=green!40,
							font=\ttfamily}}
% بسته‌ و دستوراتی برای ایجاد لینک‌های رنگی با امکان جهش
\usepackage[linktocpage=true,colorlinks,pagebackref=true,linkcolor=blue,citecolor=magenta]{hyperref}
% چنانچه قصد پرینت گرفتن نوشته خود را دارید، خط بالا را غیرفعال و  از دستور زیر استفاده کنید چون در صورت استفاده از دستور زیر‌‌، 
% لینک‌ها به رنگ سیاه ظاهر خواهند شد که برای پرینت گرفتن، مناسب‌تر است
%\usepackage[pagebackref=false]{hyperref}
%بسته‌ی لازم برای مراجع
\usepackage[nonamebreak,square]{natbib}
% بسته‌ لازم برای تنظیم سربرگ‌ها
\usepackage{fancyhdr}
% بسته‌ای برای ظاهر شدن «مراجع» و «نمایه» در فهرست مطالب
%\usepackage[nottoc]{tocbibind}
% دستورات مربوط به ایجاد نمایه
\usepackage{makeidx}
\makeindex
%%%%%%%%%%%%%%%%%%%%%%%%%%
% فراخوانی بسته زی‌پرشین و تعریف قلم فارسی و انگلیسی
\usepackage{xepersian}
\settextfont[Scale=1.15]{XB Niloofar}
%\setiranicfont[Scale=1]{PersianModern-Oblique}
\setiranicfont[Scale=1.15]{XB Zar Oblique}
% چنانچه می‌خواهید اعداد در فرمول‌ها، انگلیسی باشد، خط زیر را غیرفعال کنید
%\setdigitfont{Yas}
\setdigitfont[Scale=.85]{Persian Modern}
%\setdigitfont[Scale=1.1]{XB Zar}
%\KashidaOn
%تنظیم فاصله‌ی خطوط و پاراگراف‌ها
\linespread{1.3}
\setlength{\parskip}{1ex plus 0.5ex minus 0.2ex}
%%%%%%%%%%%%%%%%%%%%%%%%%%
% تعریف قلم‌های فارسی و انگلیسی اضافی برای استفاده در بعضی از قسمت‌های متن
\defpersianfont\nastaliq[Scale=2.02]{IranNastaliq}
\defpersianfont\chapternumber[Scale=3.02]{XB Niloofar}
\defpersianfont\titr[Scale=1.02]{XB Titre}
\defpersianfont\dav[Scale=2]{B Davat}
%%%%%%%%%%%%%%%%%%%%%%%%%%
% دستوری برای حذف کلمه «چکیده»
\renewcommand{\abstractname}{}
% دستوری برای حذف کلمه «abstract»
\renewcommand{\latinabstract}{}
% دستوری برای تغییر نام کلمه «اثبات» به «برهان»
\renewcommand\proofname{\textbf{برهان}}
% دستوری برای تغییر نام کلمه «کتاب‌نامه» به «مراجع»
\renewcommand{\bibname}{مراجع}
%استفاده از کلمه‌ی نمودار به جای شکل
\renewcommand{\figurename}{نمودار}
% دستوری برای تعریف واژه‌نامه انگلیسی به فارسی
\newcommand\persiangloss[2]{#1\dotfill\lr{#2}\\}
% دستوری برای تعریف واژه‌نامه فارسی به انگلیسی 
\newcommand\englishgloss[2]{#2\dotfill\lr{#1}\\}
%%%%%%%%%%%%%%%%%%%%%%%%%%
% تعریف و نحوه ظاهر شدن عنوان قضیه‌ها، تعریف‌ها، مثال‌ها و ...
\theoremstyle{definition}
\newtheorem{definition}{تعریف}[section]
%\newtheorem{deduct}[definition]{}
\newtheorem{theorem}[definition]{قضیه}
\newtheorem{lemma}[definition]{لم}
\newtheorem{proposition}[definition]{گزاره}
\newtheorem{corollary}[definition]{نتیجه}
\newtheorem{remark}[definition]{ملاحظه}
\newtheorem{example}[definition]{مثال}
%%%%%%%%%%%%%%%%%%%%%%%%%%
% تعریف دستورات جدید برای خلاصه نویسی و راحتی کار در هنگام تایپ فرمول‌های ریاضی
\newcommand{\mA}{\mathcal{A}}% مجموعه‌ی عامل‌ها
\newcommand{\mB}{\mathcal{B}}% زیرمجموعه‌ای از عامل‌ها
\newcommand{\mL}{\mathcal{L}}% زبان
\newcommand{\mF}{\mathcal{F}}% سیگما جبر
\newcommand{\mP}{\mathcal{P}}% تخصیص احتمالاتی
\newcommand{\mK}{\mathcal{K}}% منطق K
\newcommand{\mS}{\mathcal{S}}% منطق S5
\newcommand{\mbbP}{\mathbb{P}}% مجموعه‌ی گزاره‌های اتمی
\newcommand{\mbbQ}{\mathbb{Q}}% مجموعه‌ی اعداد گویا
\newcommand{\mbbM}{\mathbb{M}}% مجموعه‌ی مدل‌های شناختی احتمالاتی
\newcommand{\mbbK}{\mathbb{K}}% کلاس مدل‌های کریپکی
\newcommand{\mbbS}{\mathbb{S}}% کلاس مدل‌های شناختی
\newcommand{\mbP}{\mathbf{P}}% نماد تابعی احتمالاتی
\newcommand{\xra}{\xrightarrow}
\newcommand{\scr}{\scriptscriptstyle}
\newcommand{\bp}{\begin{proof}}
\newcommand{\ep}{\end{proof}}
%\newcommand{\close}{\begin{latin}\noindent $\square$ \end{latin}}
%%%%%%%%%%%%%%%%%%%%%%%%%%%%%%%%%%%%%%%%%%%
%تعریف اعداد لاتین برای زمانی که از اعداد فارسی در متن استفاده می‌کنیم
\def\0{\textrm{\lr{0}}}
\def\1{\textrm{\lr{1}}}
\def\2{\textrm{\lr{2}}}
\def\3{\textrm{\lr{3}}}
\def\4{\textrm{\lr{4}}}
\def\5{\textrm{\lr{5}}}
%%%%%%%%%%%%%%%%%%%%%%%%%%%%%%%%%%%%%%%%%%%
\def\reduction{تحویل}
\def\observation{مشاهده‌محور}
\def\prior{پیشینی}
%%%%%%%%%%%%%%%%%%%%%%%%%%%%%%%%%%%%%%%%%%%
% براساس نسخه‌های از ۱.۲.۵ به بعد بسته‌ی bidi در توابع زیر باید توجه داشت:
% که l یعنی شروع نوشتار از ابتدای خط یا جعبه،
% و r یعنی شروع نوشتار از انتهای خط یا جعبه، 
% و s یعنی شروع نوشتار از وسط خط یا جعبه.
\newdimen\xleftright
\xleftright=\textwidth
\advance \xleftright by -10.5cm
\newcommand{\leftright}[3]{%
\noindent
\makebox[4 cm][r]{#1}
\makebox[\xleftright][s]{}
\makebox[5.5 cm][l]{#2}
\makebox[1 cm][l]{#3}%
}
\newcommand{\leftrightb}[2]{%
\noindent
\makebox[4 cm][r]{#1}
\makebox[\xleftright][s]{}
\makebox[6.5 cm][l]{#2}%
}
\newcommand{\semanticsa}[2]{%
\noindent
\makebox[4.3 cm][r]{#1}
\makebox[\xleftright][c]{اگر و فقط اگر}
\makebox[6.2 cm][l]{#2}%
}
\newcommand{\semanticsb}[2]{%
\noindent
\makebox[3.5 cm][l]{#1}
\makebox[\xleftright][c]{اگر و فقط اگر}
\makebox[7 cm][r]{#2}%
}
%%%%%%%%%%%%%%%%%%%%%%%%%%%%
% دستورهایی برای سفارشی کردن سربرگ صفحات
\csname@twosidetrue\endcsname
\pagestyle{fancy}
\fancyhf{} 
\fancyhead[RE,LO]{\thepage}
\fancyhead[LE]{\small\iranicfamily\leftmark}
\fancyhead[RO]{\small\iranicfamily\rightmark}
\renewcommand{\chaptermark}[1]{%
\markboth{\thechapter.\ #1}{}}
%%%%%%%%%%%%%%%%%%%%%%%%%%%%%
% دستورهایی برای سفارشی کردن صفحات اول فصل‌ها
\makeatletter
\newcommand\mycustomraggedright{%
 \if@RTL\raggedleft%
 \else\raggedright%
 \fi}
\def\@part[#1]#2{%
\ifnum \c@secnumdepth >-2\relax
\refstepcounter{part}%
\addcontentsline{toc}{part}{\thepart\hspace{1em}#1}%
\else
\addcontentsline{toc}{part}{#1}%
\fi
\markboth{}{}%
{\centering
\interlinepenalty \@M
\ifnum \c@secnumdepth >-2\relax
 \huge\bfseries \partname\nobreakspace\thepart
\par
\vskip 20\p@
\fi
\Huge\bfseries #2\par}%
\@endpart}
\def\@makechapterhead#1{%
\vspace*{-30\p@}%
{\parindent \z@ \mycustomraggedright %\@mycustomfont
\ifnum \c@secnumdepth >\m@ne
\if@mainmatter

\huge\bfseries \@chapapp\space {\chapternumber\thechapter}
\par\nobreak
\vskip 20\p@
\fi
\fi
\interlinepenalty\@M 
\Huge \bfseries #1\par\nobreak
\vskip 120\p@
}}
\makeatother
\begin{document}
\frontmatter
%%%%%%%%%%%%%%%%%%%%%%%%%%%%%%%%%%%%%%%%%
%صفحه‌ی بسم الله
\thispagestyle{empty}
\centerline{{\includegraphics[scale=2.5]{besmelah.jpg}}}
\newpage
\thispagestyle{empty}
\clearpage
%~~~
%%%%%%%%%%%%%%%%%%%%%%%%%%%%%%%%%%%%%%%%%
\clearpage
\pagenumbering{adadi}
% در این فایل، عنوان پایان‌نامه، مشخصات خود، متن تقدیمی‌، ستایش، سپاس‌گزاری و چکیده پایان‌نامه را به فارسی، وارد کنید.
% توجه داشته باشید که جدول حاوی مشخصات پایان‌نامه/رساله و همچنین، مشخصات داخل آن، به طور خودکار، درج می‌شود.
%%%%%%%%%%%%%%%%%%%%%%%%%%%%%%%%%%%%
% دانشگاه خود را وارد کنید
\university{ شهید بهشتی}
% دانشکده، آموزشکده و یا پژوهشکده  خود را وارد کنید
\faculty{علوم ریاضی}
% گروه آموزشی خود را وارد کنید
\degree {کارشناسی ارشد} 
% گروه آموزشی خود را وارد کنید
\subject{علوم‌کامپیوتر }
% گرایش خود را وارد کنید
\field{علوم‌کامپیوتر}
% عنوان پایان‌نامه را وارد کنید
\title{کنترل همروندی در پایگاه داده های ذاتا XML}
% نام استاد(ان) راهنما را وارد کنید
\firstsupervisor{دکتر سید کامیار ایزدی}
%\secondsupervisor{استاد راهنمای دوم}
% نام استاد(دان) مشاور را وارد کنید. چنانچه استاد مشاور ندارید، دستور پایین را غیرفعال کنید.
\firstadvisor{دکتر محمود فضلعلی}
%\secondadvisor{استاد مشاور دوم}
% نام پژوهشگر را وارد کنید
\name{محسن }
% نام خانوادگی پژوهشگر را وارد کنید
\surname{لچینانی}
% تاریخ پایان‌نامه را وارد کنید
\thesisdate{1394}
% کلمات کلیدی پایان‌نامه را وارد کنید
\keywords{کنترل همروندی ، Native XML، Query Processing}
% چکیده پایان‌نامه را وارد کنید
\fa-abstract{\noindent
در این پایان‌نامه ابتدا نمونه‌ای از منطق‌های شناختی احتمالاتی (PEL) را معرفی کرده و تمامیت آن را اثبات می‌کنیم. سپس با در نظر گرفتن مدلی ساده‌ از این منطق، در راستای تعمیم آن به منطق‌های پویا که تغییر اطلاعات در سناریوهای چند عاملی را مدل می‌کنند پیش می‌رویم.
\\
پس از توصیف مختصری از منطق‌های شناختی پویای غیر احتمالاتی، منطق شناختی پویای احتمالاتی (PDEL) را نیز با در نظر گرفتن سه گونه‌ی طبیعی احتمال، یعنی احتمال {\prior} جهان‌ها، احتمال رخداد عمل‌ها بر اساس فرایندی متناظر با دیدگاه عامل‌ها و احتمال خطا در مشاهده‌ی عمل‌ها، معرفی خواهیم کرد. این سه گونه، شیوه‌ی به‌روزرسانی تعمیم‌یافته‌ای در اختیار می‌گذارند که روشی است مناسب و طبیعی برای مدل‌سازی جریان اطلاعات.
\\
سپس برای اینکه تمامیت منطق شناختی پویای احتمالاتی را با استفاده از تمامیت منطق شناختی احتمالاتی اثبات کنیم، اصول موضوعه‌ی صحیحی ارایه می‌کنیم تا فرمول‌های شامل عملگر پویا را به فرمول‌هایی فاقد این عملگر در زبان ایستای متناظر تحویل کنند.
\\
سرانجام گونه‌ای از منطق شناختی پویای احتمالاتی، مورد نیاز برای حل معمای \lr{Monty  Hall}\!\!\! ، ارایه کرده و تمامیت آن را اثبات می‌کنیم. سپس راه‌حلی صوری برای این معما در این منطق بدست می‌آوریم. 
}
\newpage
\thispagestyle{empty}
\vtitle
\newpage
\thispagestyle{empty}
\clearpage
~~~
\newpage
\thispagestyle{empty}
\ \\ \\ \\ \\ \\ \\ \\
{\dav
\begin{center}
كلية حقوق اعم از چاپ و تكثير، نسخه برداري ، ترجمه، اقتباس و ... از اين پايان نامه براي دانشگاه شهيد بهشتي محفوظ است.
 نقل  مطالب با ذكر مأخذ آزاد است.
\end{center}
}
\newpage
\thispagestyle{empty}
\clearpage
~~~
%\newpage
%\thispagestyle{empty}
%\centerline{{\includegraphics[width=20 cm]{replyrecord}}}
%\newpage
%\thispagestyle{empty}
%\clearpage
%~~~
\newpage
 % پایان‌نامه خود را تقدیم کنید!
\begin{acknowledgementpage}

\vspace{4cm}

{\nastaliq
{\Large
تقدیم به پدر و مادر عزیزم که همواره در تمامی مراحل زندگی پشتیبان و حامی من بوده‌اند.

\vspace{1.5cm}

\newdimen\xa
\xa=\textwidth
\advance \xa by -11cm
\hspace{\xa}

}}
\end{acknowledgementpage}
\newpage
\thispagestyle{empty}
\clearpage
~~~
%%%%%%%%%%%%%%%%%%%%%%%%%%%%%%%%%%%%
\newpage
\thispagestyle{empty}
% ستایش
\baselineskip=.750cm
\ \\ \\
 
{\nastaliq
رسيدن، به دانش است و به كردار نیک...%
}\\
\vspace{.5cm}\\
{\scriptsize\nastaliq
{
 و بی دانش به كردار نيك هم نتوان رسيد، كه نيكى را پيشتر ببايد شناختن، آنگاه بجاى آوردن. پس دانش به همه حال مى‌ببايد تا به رستگارى توان رسيدن. و چون دانش راه آمد، به بهترين چيزها كه آدمى را تواند بودن. و در اوّل آفرينش حاصل نيست و بعضى از آن بى‌رنج و انديشه حاصل شود، پس هرآينه مهمتر چيزى باشد كه در حاصل كردنش عمر گذرانند، ليكن برخى هست كه بى‌انديشه حاصل آيد و بعضى را ناچار به انديشه حاجت بود، و آنچه به انديشه حاصل شود دانسته‌اى خواهد كه درو انديشه كنند تا اين نادانسته بدان انديشه كه در آن دانسته كنند دانسته شود، و از هر دانسته هر نادانسته را نتوان شناخت، بلكه هر نادانسته را به دانسته‌اى كه در خور او بود توان شناخت. و منطق آن علم است كه درو راه انداختن نادانسته به دانسته دانسته شود...
 }}
 
\vspace{.5cm}
{\nastaliq
\newdimen\xb
\xb=\textwidth
\advance \xb by -8.5cm
\hspace{\xb}
پس منطق ناگزير آمد بر جوينده‌ی رستگاری.
\RTLfootnote{مقدمه‌ی رساله‌ی منطق دانشنامه‌ی علائی، شیخ‌الرئیس ابن‌سینا}
}
\newpage
\thispagestyle{empty}
\clearpage
~~~
%%%%%%%%%%%%%%%%%%%%%%%%%%%%%%%%%%%%
\newpage
\thispagestyle{empty}
% سپاس‌گزاری
{\nastaliq
سپاس‌گزاری...
}
\\[2cm]
ستایش و سپاس مخصوص خداوندی ست که مرا در مسیر علم و دانش قرار داد.

ارزشمندترین زمان برای نیک‌اندیشان، لحظه‌ای است که در محضر استاد می نشینند و از خرمن فضلش خوشه می‌چینند. بسیار مایلم صمیمانه‌ترین قدردانی و سپاس خویش را تقدیم استاد گرامی، جناب آقای دکتر ایزدی سازم که در طول این دوره توفیق راهم شدند تا در حلقه شاگردیشان قرار گیرم و از مرتبه‌های علم و دانش ایشان بهره‌ها ببرم. همچنین از سركار خانم دکتر طهماسبی که راهنمایی‌های ایشان در انجام این پروهش همواره چراغ راهم بوده است. 



% با استفاده از دستور زیر، امضای شما، به طور خودکار، درج می‌شود
\signature 
\newpage
\thispagestyle{empty}
\clearpage
~~~
\newpage
%{\small
\abstractview
%}
\newpage
\thispagestyle{empty}
\clearpage
~~~
\newpage
\chapter*{پیش‌گفتار}
\markboth{پیش‌گفتار}{پیش‌گفتار}

منطق شناختی احتمالاتی پویا منطقی نسبتاً جدید است که قبل از مطالعه‌ی آن می‌بایست منطق شناختی، منطق شناختی پویا و منطق شناختی احتمالاتی معرفی شده باشد و به فراخور در فصول سه‌گانه‌ به معرفی و توصیف هریک خواهیم پرداخت. 

مقدمه را با بهره‌جستن از مقدمه‌ی مقاله‌ی \citep{Kooi2003} نگاشتم و البته هر کجا که لازم بود از مقدمه‌های دیگر مقالات نیز استفاده کردم.

در فصل 1 ابتدا با کمک \citep{DELDitmarsch2007} و \citep{Sack2007} به خلاصه‌ای از منطق شناختی می‌پردازیم، سپس بر مبنای \citep{Fagin1994} منطق شناختی احتمالاتی را که همان منطق شناختی سنتی است به اضافه‌ی توانایی استدلال درباره‌ی احتمال معرفی کرده و تمامیت آن را اثبات می‌کنیم. 

فصل 2 بر اساس \citep{Benthem2009} نوشته شده است و به معرفی منطق‌های پویا اعم از شناختی پویا و شناختی پویای احتمالاتی اختصاص دارد. در این فصل پس از معرفی منطق اعلان عمومی و سپس گونه‌ای تا حدی تعمیم یافته از منطق شناختی پویا با در نظر گرفتن مدلی ساده‌ از منطق‌ شناختی احتمالاتی، منطق شناختی پویای احتمالاتی معرفی شده و تمامیت آن اثبات می‌شود. برای آنکه منطق‌های شناختی پویا را احتمالاتی کنیم ابتدا سه گونه‌ی طبیعی احتمال را تعریف می‌کنیم که عبارتند از: احتمال {\prior} جهان‌ها، احتمال رخداد عمل‌ها بر اساس فرایندی متناظر با دیدگاه عامل‌ها و احتمال خطا در مشاهده‌ی عمل‌ها. برای اثبات تمامیت منطق‌های پویا اصول موضوعه‌ای مطرح می‌شود تا بتوان معادل با هر فرمول در این منطق‌ها، با حذف عملگر پویا، فرمولی در منطق‌های ایستا بدست آورد. این اصول اثرات متقابل عملگر پویا با اتم‌ها و عملگرهای بولی و شناختی را توصیف می‌کنند. پس از اثبات صحت، با کمک آنها تمامیت منطق‌های پویا را از تمامیت منطق‌های ایستا نتیجه می‌گیریم.

هنگامی که به همراه ولی‌زاده به مطالعه و اثبات جزئیات مقاله‌ی \citep{Benthem2009} مشغول بودیم با مسائل و مشکلاتی برخورد کردیم که به تناسب در بخش‌های مختلف فصل ۲ با عنوان ملاحظه به آنها اشاره خواهم کرد.

فصل 3 نیز بر گرفته از \citep{Kooi2003} است و به منظور ارائه‌ی مثالی برای نشان دادن کاربرد منطق شناختی پویای احتمالاتی نگاشته شده است. مثالی که ارائه می‌شود معمای معروف \lr{Monty Hall} است که پیش از این مقاله راه‌حلی صوری برای آن داده نشده بود، ولی ما بر اساس این مقاله به راه‌حلی صوری با کمک گونه‌ای از منطق‌های شناختی پویای احتمالاتی دست می‌یابیم که به زیبایی جوابی معقول در اختیار می‌گذارد. این راه‌حل زیبا را نمی‌توانستم با جزئیات بیان کنم اگر از \citep{Kooithesise2003} استفاده نمی‌کردم و آن را نیافتم مگر با راهنمایی دکتر کویی.

\tableofcontents
%\listoffigures
\clearpage{\pagestyle{empty}\cleardoublepage}
%برای توضیح این دستور به راهنمای بسته‌ی fancyhdr مراجعه شود
\mainmatter
\clearpage
\phantomsection
\addcontentsline{toc}{chapter}{مقدمه}
\chapter*{مقدمه}\markboth{مقدمه}{مقدمه}

در سال ۱۹۵۱ فان رایت\LTRfootnote{\lr{von Wright}} کتابی با عنوان «مقاله‌ای در منطق موجهات» \citep{Wright1951} منتشر کرد. هینتیکا\LTRfootnote{\lr{Hintikka}} با ایده‌هایی که از این مقاله گرفته بود در سال ۱۹۶۲ کتابی با عنوان «دانش و باور، مقدمه‌ای بر منطق مبتنی بر این دو مفهوم» \citep{Hintikka1962} به چاپ رساند. وی در این کتاب به کمک مفهوم جهان‌های ممکن مدلی برای دانش و باور ارائه کرد و به همین دلیل بسیاری او را پدر منطق شناختی می‌دانند. هدف اصلی او واکاوی مفهومی دانش و باور بود ولی پس از او عبارت «شناخت» در محدوده‌ای فراتر به کار گرفته شد، اعم از باور و هر روشی که یک عامل می‌تواند دانشی را بدست آورد. در اواخر دهه‌ی ۱۹۷۰ منطق شناختی مورد توجه دانشمندان فعال در شاخه‌هایی مانند هوش مصنوعی، فلسفه و نظریه بازی‌ها قرار گرفت.
در دهه ۸۰ محققین علوم کامپیوتر به منطق شناختی روی آوردند، فاگین\LTRfootnote{\lr{Fagin}}، هالپرن\LTRfootnote{\lr{Halpern}}، موزِز\LTRfootnote{\lr{Moses}} و وَردی\LTRfootnote{\lr{Vardi}} که از این دسته به حساب می‌آمدند مقالاتی را که در طی حدود ۱۰ سال در مورد منطق شناختی به چاپ رسانده بودند در کتابی با نام «استدلال درباره دانش» \citep{RAKFagin1995} جمع آوری کرده و در سال ۱۹۹۵ منتشر کردند.

منطق شناختی، منطقی وجهی است که به استدلال برمبنای دانش\LTRfootnote{\lr{information}} و فرادانش\LTRfootnote{\lr{higher order information}}\index{فرادانش} می‌پردازد، فرادانش دانشی است که عاملی درباره‌ی دانش خود و یا دانش دیگر عامل‌ها دارد. برای روشن شدن موضوع مثالی را با سه بازیکن 1، 2 و 3، و سه کارت قرمز، سفید وآبی در نظر بگیرید. کارت‌ها در میان بازیکنان به این صورت توزیع شده‌اند که قرمز، سفید و آبی به ترتیب در دستان 1، 2 و 3 قرار دارد. فرض کنید بازیکنان تنها کارت خویش را می‌بینند و همه می‌دانند که کارت‌ها به گونه‌ای میانشان توزبع شده است که هریک فقط یک کارت در دست دارند.  بوسیله‌ی منطق شناختی می‌توان جملات پیچیده‌ای را مانند «بازیکن 1 می‌داند که بازیکن 2 نمی‌داند که چه کارتی در دستان بازیکن 3 است.» صوری کرد. منطق شناخت با چنین فرادانش‌هایی سر و کار دارد. حتی فرادانش‌های پیچیده‌تری مانند همه‌دانی مشترک وجود دارد. برای نمونه در مثال بالا علم به اینکه دقیقاً سه کارت وجود دارد و علم به رنگ کارت‌ها همه‌دانی مشترک است. منطق شناختی همچنین واکاوی مفهومی خوبی از همه‌دانی مشترک در اختیار می‌گذارد.

اگرچه منطق شناختی آنالیز مناسبی از فرادانش در اختیار می‌گذارد ولی بررسی تغییر دانش خارج از گستره‌ی این منطق است. منطق‌های شناختی پویا منطق‌ شناختی را به گونه‌ای توسیع می‌دهند که استدلال درباره‌ی تغییرات دانش نیز امکان‌پذیر باشد. این توسیع از طرفی از نوعی معناشناسی زبان طبیعی الهام گرفته شد که در آن معنای جمله بعنوان طریقی برای تغییر داده‌های کسانی که آن را می‌شنوند در نظر گرفته می‌شود، و از طرف دیگر از مطالعه‌ی بازی‌ها که تغییر داده‌ها و فرادانش‌ها نقش بسزایی در آنها ایفا می‌کند. سیستم‌های منطقی مختلفی بر این اساس در طول سال‌ها شکل گرفته است که برجسته‌ترین آنها عبارتند از \citep{Gerbrandy1997} (الهام گرفته از \citep{Veltman1996})، \citep{Batlag1998} و  \citep{Cate2002}.

در منطق شناختی پویا تغییر وضعیت موجود توسط داده‌های جدید را به‌روزرسانی می‌خوانیم. ساده‌ترین مثال زمانی است که عاملی می‌فهمد که گزاره‌ی $ \varphi $ برقرار است. به‌روزرسانی با یک گزاره در این مثال به این معنی است که گزینه‌هایی که عامل ممکن می‌دانست ولی در آنها گزاره‌ی $ \varphi $ برقرار نیست حذف می‌شوند. در یک سیستم چند عاملی ممکن است عامل‌های مختلف دسترسی مختلفی به داده‌های جدید داشته باشند و همچنین اطلاعات عامل‌ها درباره‌ی دیگر عامل‌ها نقش بازی کند، از این رو می‌توان به‌روزرسانی پیچیده‌تری را در مثال قبل مدل کرد: فرض کنید بازیکن 1 کارت خود را به بازیکن 2 نشان دهد و بازیکن 3 نیز این را ببیند ولی از محتوای کارت خبردار نشود. در نتیجه دانش بازیکنان به این صورت تغییر می‌کند: بازیکن 2 می‌داند که محتوای کارت بازیکن 1 چیست، بازیکن 3 می‌داند که بازیکن 2 می‌داند محتوای کارت بازیکن 1 جیست ولی خودش محتوا را نمی‌داند و بازیکن 1 می‌داند که بازیکن 3 این را می‌داند.

اگرچه نظریه‌ی احتمال منطق نیست لکن حوزه‌ی مطالعاتی مناسبی برای منطق است، زیرا در بسیاری از حوزه‌های کاربردی به منظور استدلال درباره‌ی دانش، اهمیت توانایی استدلال درباره‌ی احتمال رخدادهای معین به همراه دانش عامل‌ها رخ می‌نماید و اغلب احتمال به عنوان نظریه‌ای برای مدلسازی استدلال مطرح می‌شود.
از این رو همه‌ی مقالات منتشر شده در علم اقتصاد که به استدلال درباره‌ی دانش می‌پردازند (که بازگشت می‌کنند به مقاله‌ی اصلی اومان \citep{Aumann1976}) با ساختاری احتمالاتی مدل می‌شوند، هرچند آنها زبانی منطقی که بصورتی روشن استدلال درباره‌ی احتمال را جایز کند در نظر نگرفته‌اند. با این اوصاف تلاش‌هایی در جهت بیرون کشیدن منطق بعنوان بهترین راه استدلال از دل نظریه‌ی احتمال صورت گرفته است. که یکی از مناسب‌تزین آنها منطق احتمالاتی است که در \citep{Fagin1990} معرفی شده است.
\clearpage{\pagestyle{empty}\cleardoublepage}
\chapter{منطق‌های شناختی ایستا}
منطق شناختی چیز خوبی است.
\section{منطق شناختی \texorpdfstring{ \lr{(EL)}}{(EL)}}\index{منطق!شناختی}

\clearpage{\pagestyle{empty}\cleardoublepage}
\chapter{منطق‌های شناختی پویا}
در این فصل سه منطق شناختی پویا را با این رویکرد مطرح می‌کنیم که برای هریک اصولی موسوم به اصول موضوعه‌ی {\reduction\LTRfootnote{\lr{reduction axioms}}} معرفی کرده و با اثبات صحت آنها گامی به سوی تمامیت بر می‌داریم. در انتهای فصل نیز تمامیت را در یک قضیه برای هر سه منطق اثبات خواهیم کرد.
\section{منطق‌های شناختی پویا به منظور به‌روزرسانی غیر احتمالاتی}
منطق‌های شناختی پویا جریان اطلاعات ایجاد شده توسط عمل\LTRfootnote{\lr{event}}‌ها را توصیف می‌کنند. ساده‌ترین عمل آموزنده، و نمونه‌ای رهگشا برای بیشتر این نظریه، اعلان عمومیِ\index{اعلان عمومی} گزاره‌یِ درستی چون $ A $ به گروهی از عامل‌هاست، که به‌صورت $ !A $  نمایش می‌دهیم. به‌روزرسانی برای عمل‌های پیچیده‌تر می‌تواند برحسب «مدل‌های عمل» توصیف شود، که الگوهای پیچیده‌تری از دسترسی عامل‌ها به عملِ در حال رخداد را مدل می‌کنند. پس ابتدا به‌روزرسانی منطق شناختی توسط اعلان عمومی را بررسی می‌کنیم سپس آن را به حالت کلی‌تر، برای هر نوع عمل، توسیع می‌دهیم.

\subsection{منطق اعلان عمومی \texorpdfstring{ \lr{(PAL)}}{(PAL)}}\index{منطق!اعلان عمومی}
تأثیر پویای اعلان عمومی\LTRfootnote{\lr{public announcement}} $ A $ این است که مدل (غیر احتمالاتی) جاری  $ M=(S,\sim,V) $ را به مدل به‌روز شده‌ی $ M|A $ تبدیل می‌کند. این مدل به‌روز شده با تحدید جهان‌های $ M $ به جهان‌هایی که $ A $ در آنها درست است تعریف می‌شود.

اعلان عمومی معمولاً حاوی اطلاعاتی مفید است. از این رو ممکن است که ارزش درستی عبارات شناختی در نتیجه‌ی اعلان تغییر کند. برای مثال قبل از اعلان $ A $ عامل $ a $ آن را نمی‌دانست ولی اکنون می‌داند.

\begin{definition}{\textbf{زبان اعلان عمومی.}}\index{زبان!اعلان عمومی}
زبان اعلان عمومی توسط فرم \lr{Backus-Naur} به‌صورت زیر بیان می‌شود:
\begin{equation*}
\varphi,\psi ::=\ \top\mid\bot\mid p \mid\neg\varphi\mid\varphi\wedge\psi\mid K_i\varphi\mid\left[ !\varphi\right] \psi
\end{equation*}
\end{definition}
فرمول $ [!\varphi]\psi $ به‌صورت «$ \psi $ پس از اعلان $ \varphi $ برقرار است» خوانده می‌شود. زبان بدست آمده  در مدل‌های استاندارد برای منطق شناختی نیز قابل تفسیر است. معناشناسی برای این زبان به غیر از اعلان عمومی همانند تعریف \ref{def3} می‌باشد. معناشناسی اعلان عمومی نیز به‌صورت زیر تعریف می‌شود.
\begin{definition}{\textbf{معناشناسی اعلان عمومی.}}\index{معناشناسی!اعلان عمومی}
فرض کنید مدل شناختی  $ M=(S,\sim,V) $ داده شده باشد و $ s\in S $.
\\

\semanticsb{$ M,s\vDash [!A]\varphi $}{اگر $ M,s\vDash A $ آنگاه $ M|A,s\vDash\varphi $}
\\
\\
که در آن $ M|A $ مدل $ (S',\sim ',V') $ است به طوری که، با فرض
$ \llfloor A\rrfloor =\{t\in S\mid M,t\vDash A\} $:
\begin{LTR}
\begin{itemize}
\item
$ S'=\llfloor A\rrfloor, $
\item
$ \sim'_a=\sim_a\cap(S'\times S'), $
\item
$ V'(p)=V(p)\cap S'. $
\end{itemize}
\end{LTR}
\end{definition}
اصول موضوعه‌‌ی {\reduction} در PAL به‌صورت زیر است:
\begin{align}
&[!A]p\leftrightarrow(A\rightarrow p)\label{1}\\
&[!A]\neg\varphi\leftrightarrow(A\rightarrow\neg[!A]\varphi)\label{2}\\
&[!A](\varphi\wedge\psi)\leftrightarrow([!A]\varphi\wedge[!A]\psi)\label{3}\\
&[!A]K_a \varphi\leftrightarrow(A\rightarrow K_a[!A]\varphi)\label{4}
\end{align}

\begin{theorem}\label{reduct1}\textbf{(صحت اصول موضوعه‌ی  {\reduction} برای اعلان عمومی)}
\end{theorem}
\bp
با ارجاع به هر اصل اثباتی برای آن می‌آوریم.
\begin{itemize}
\item[(\ref{1})]
\begin{align*}
M,s\vDash [!A]p &\ \ \Leftrightarrow\ \ M,s\vDash A\Rightarrow M|A,s\vDash p\tag{1}\\
&\ \ \Leftrightarrow\ \ M,s\vDash A\Rightarrow M,s\vDash p\tag{2}\\
&\ \ \Leftrightarrow\ \ M,s\vDash A\rightarrow p
\end{align*}
اگر $ M,s\vDash A $ آنگاه $ s\in S' $ و اگر $ s\in S' $  آنگاه $ V(p)=V'(p) $. در نتیجه از (1) به (2) و برعکس می‌توان رسید.
\item[(\ref{2})]
\begin{align*}
M,s\vDash[!A]\neg\varphi &\ \ \Leftrightarrow\ \ M,s\vDash A\Rightarrow M|A,s\vDash\neg \varphi\\
&\ \ \Leftrightarrow\ \ M,s\vDash A\Rightarrow (M,s\vDash A\ \textrm{و}\  M|A,s\nvDash\varphi)\\
&\ \ \Leftrightarrow\ \ M,s\vDash A\Rightarrow M,s\vDash\neg[!A]\varphi\\
&\ \ \Leftrightarrow\ \ M,s\vDash A\rightarrow\neg[!A]\varphi
\end{align*}
\end{itemize}
\ep
\subsection{منطق شناختی پویا - به‌روزرسانی مدل‌ها \texorpdfstring{ \lr{(DEL)}}{(DEL)}}\index{منطق!شناختی!پویا}
\begin{definition}{\textbf{مدل عمل\LTRfootnote{\lr{event model}}.}}\index{مدل!عمل}
فرض کنید مجموعه‌ی $ \mA $ از عامل‌ها و زبان منطقی $ \mL $ داده شده باشد، مدل عمل ساختار $ A=(E,\sim,pre) $ است بطوری که
\begin{itemize}
\item
$ E $ مجموعه‌ای متناهی و غیر تهی است از عمل‌ها،
\item
$ \sim $ مجموعه‌ای است از روابط هم‌ارزی $ \sim_a $ روی $ E $ برای هر عامل $ a\in\mA $،
\item
$ pre $ تابعی است که به هر عمل $ e\in E $ فرمولی از $ \mL $ را نسبت می‌دهد.
\end{itemize}
تابع پیش‌شرطِ\LTRfootnote{\lr{precondition function}}\index{تابع پیش‌شرط} $ pre $ با نسبت دادن فرمول $ (pre_e) $ به هر عمل در $ E $ معین می‌کند که در کدام جهان‌ها این عمل‌ها ممکن است روی دهند. این مدل‌ها را مدل به‌روزرسانی نیز می‌نامند.
\end{definition}
این مدل‌ها بسیار شبیه مدل‌های شناختی هستند، با این تفاوت که به‌جای دانش‌های مربوط به وضعیت‌های ثابت، دانش درباره‌ی عمل‌ها \index{عمل}مدل شده است.\RTLfootnote{کلمه‌ی «عمل» ترجمه‌ای است از کلمه‌ی \lr{event}، از آنجایی که این کلمه علاوه بر منطق شناختی پویا در نظریه احتمالات نیز استفاده می‌شود، باید دانست که با تفسیرهای متفاوتی در این دو مقوله به کار می‌رود. در نظریه احتمال، \lr{event} آن است که در منطق بدان گوییم گزاره. در حالی که یک \lr{event} در منطق شناختی پویا به همراه گزاره‌ی پیش‌شرط ایجاد می‌شود، ولی در واقع \lr{event}های مدلِ عمل، مدلِ شناختی داده شده را تغییر می‌دهند و خود بخشی از مدل نیستند. از این پیچیده‌تر، گاهی اوقات به تمام مدل عمل، یک \lr{event} اطلاق می‌شود.} روابط تمییز ناپذیری $ \sim $ روی عمل‌ها ابهام درباره‌ی اینکه چه عملی واقعاً رخ داده است را مدل می‌کنند. $ e\sim_a e' $ می‌تواند به این صورت خوانده شود که «اگر فرض شود که عمل $ e $ رخ داده است رخداد عمل $ e' $ با دانشِ $ a $ سازگار است». 

اعلان عمومی $ [!\varphi] $ نیز خود به نوعی یک مدل عمل است که در آن $ E=\{!\} $ و  $ \sim=\{(!,!)\}  $ و $ pre=\{(!,\varphi)\} $.

نتیجه‌ی رخداد یک عمل نمایش داده شده با $ A $ در وضعیت نمایش داده شده با $ M $ برحسب ساختاری ضربی مدل می‌شود.
\section{منطق شناختی پویای احتمالاتی\texorpdfstring{ \lr{(PDEL)}}{(PDEL)}}\index{منطق!شناختی!پویا!احتمالاتی}
برای اینکه بتوانیم به گونه‌ای صریح و شفاف در باب تغییر داده‌های احتمالاتی در قالبی شناختی-پویا استدلال کنیم، می‌بایست منطق شناختی احتمالاتی موجود را به‌وسیله‌ی اصول موضوعه‌ی {\reduction}‌ مناسب توسعه دهیم. در این بخش نشان می‌دهیم که چگونه می‌توان این کار را بر مبنای معناشناسی مدل‌های عمل احتمالاتی، که معرفی خواهد شد، انجام داد.
\begin{definition}\textbf{زبان شناختی پویای احتمالاتی.}\index{زبان!شناختی!پویا!احتمالاتی}
زبان شناختی پویای احتمالاتی به فرم \lr{Backus-Naur} به‌صورت زیر معرفی می‌شود:

\begin{equation*}
\varphi,\psi ::=\ \top\mid\bot\mid p\mid \neg\varphi\mid\varphi\land\psi\mid K_a\varphi\mid [A,e]\varphi\mid\sum_{i=1}^n r_i \mbP_a(\varphi_i)\geq r
\end{equation*}
با همان نمادگذاری منطق شناختی احتمالاتی، علاوه بر آن $ A $ مدل عمل احتمالاتی و $ e $ عملی از آن می‌باشد. فرمول‌هایی که پیش‌شرط‌ها را در مدل احتمالاتی عمل تعریف می‌کنند از همین زبانی که معرفی شد می‌آیند.

در این زبان علاوه بر خلاصه‌نویسی‌های پیش‌گفته خلاصه‌نویسی‌های زیر نیز مطرح است:

$$\langle A,e \rangle\psi :\quad \neg[A,e]\neg\psi$$
و به منظور اینکه پیش‌شرط‌ها را در یک شئ از زبان فرموله کنیم قرار می‌دهیم
\begin{equation}\label{pg0}
pre_{A,e} :\quad\bigvee_{\varphi\in\Phi,pre(\varphi,e)>0}\varphi
\end{equation}
\end{definition}
\begin{remark}\label{preAe}
در مقاله‌ی \citep{Benthem2009}، $ pre_{A,e} $ به‌صورت زیر مطرح شده است:
\begin{equation}\label{pgeq0}
pre_{A,e} :\quad\bigvee_{\varphi\in\Phi,pre(\varphi,e)\geq 0}\varphi
\end{equation}
این تعریف معادل است با $ \bigvee_{\varphi\in\Phi}\varphi $ زیرا $ pre(\varphi,e) $ تابع احتمال است و همواره بزرگتر یا مساوی صفر است.

به دلایلی که مطرح می‌شود تعریف \ref{pg0} طبیعی‌تر به نظر می‌رسد. اولاً $ pre_{A,e} $ به‌عنوان پیش‌شرط $ e  $ مطرح است پس باید شامل پیش‌شرط‌هایی باشد که به $ e $ احتمال مثبت نسبت می‌دهند، ثانیاً اگر برای هر پیش‌شرط $ \varphi $ داشته باشیم $ pre(\varphi,e)=0 $ می‌توان $ pre_{A,e} $ را تعریف کرد $ \bot $ که از دو جنبه‌ی زیر قابل دفاع است:
\begin{itemize}
\item[-]
از منظر جبری وقتی ترتیب به‌وسیله‌ی استلزام روی فرمول‌ها تعریف شده باشد داریم $ \bot=\bigvee \phi $.
\item[-]
از نقطه نظر منطقی از آنجایی که منظور ما از $ pre(\varphi,e) $ احتمال رخداد $ e $ است وقتی $ \varphi $ برقرار است، زمانی که برای هر $ \varphi\in\Phi $ داریم $ pre(\varphi,e)=0 $، $ e $ از جهت احتمالاتی امکان وقوع ندارد، بنابراین اگر ما $ pre_{A,e} $ را قرار دهیم $ \bot $ از برقراری پیش‌شرط‌های $ e $ جلوگیری به عمل آورده‌ایم و از این رو اجازه نمی‌دهیم $ e $ رخ دهد.
\end{itemize}
\end{remark}
\clearpage{\pagestyle{empty}\cleardoublepage}
\chapter{معمای \texorpdfstring{ \lr{\textsc{Monty Hall}}}{\textsc{Monty Hall}}}
این فصل اختصاص یافته به بررسی مسأله‌ای دشوار، مشهور به معمای \lr{\textit{Monty Hall}}، با کمک منطق شناختی پویای احتمالاتی. همانطور که در بخش \ref{probabilitysource} مطرح شد تنها دو گونه از سه گونه احتمالی که معرفی شد در به‌روزرسانی در این معما کافی است و همچنین تنها عملی که در آن رخ می‌دهد اعلان عمومی است. از این رو می‌توانیم منطق را به همان دو گونه از احتمال و اعلان عمومی محدود کنیم ولی کمی از سادگی مدل‌های ایستا و پویای آن بکاهیم و بنابراین برای حل این معما یک منطق اعلان عمومی احتمالاتی را به گونه‌ای معرفی می‌کنیم که مدل‌های ایستای آن مدل‌های کریپکی احتمالاتی باشند.
\section{منطق اعلان عمومی احتمالاتی \texorpdfstring{ \lr{(PPAL)}}{(PPAL)}}
مدل‌های شناختی احتمالاتی معرفی شده در فصل 1 را به یاد بیاورید، در مدل کریپکی احتمالاتی که در اینجا معرفی می‌کنیم فضای احتمالاتی که به هر عامل $ a\in\mA $ در هر جهان $ s\in S $ تخصیص می‌دهیم به‌صورتی است که فضای نمونه‌ یعنی $ S_{a,s} $ هر زیرمجموعه‌ی دلخواهی از $ S $ می‌تواند باشد ولی $ \sigma $-جبر مجموعه‌های اندازه‌پذیر ($ \mF_{a,s} $) همواره مجموعه‌ی توانی فضای نمونه است. همچنین چون تمامی زیرمجموعه‌های تک عضوی از فضای نمونه در $ \sigma $-جبر قرار می‌گیرند، می‌توان اندازه‌ی احتمالاتی را مستقیماً روی فضای نمونه تعریف کرد و در نتیجه تعریف مدل‌های کریپکی احتمالاتی به‌صورت زیر در می‌آید:

\begin{definition}{\textbf{مدل‌های کریپکی احتمالاتی.}}\index{مدل!کریپکی!احتمالاتی}
فرض کنید $ \mA $ مجموعه‌ی عامل‌ها و $ \mbbP $ مجموعه‌ی گزاره‌های اتمی باشد. مدل کریپکی احتمالاتی ساختار $ M_{PKL}=(S,\xra{\scr{\mA}},P,V) $ است بطوریکه
\begin{itemize}
\item
$ S $ مجموعه‌ای است غیر تهی از جهان‌های ممکن،
\item
 $ \xra{\scr{\mA}} $ مجموعه‌ای است از روابط دسترسی $ \xra{a} $ که به‌ازای هر $ a\in\mA $ روی $ S $ تعریف شده‌اند،
\item
$ P:\mA\rightarrow (S\rightarrow(S_{a,s}\rightarrow[0,1])) $، تابعی احتمالاتی روی $ S_{a,s} $ به هر عامل $ a\in\mA $ و هر جهان $ s\in S $ نسبت می‌دهد (احتمالی که به $ t $ توسط تابعی که به $ a $ در $ s $ مربوط شده است نسبت داده می‌شود به‌صورت $ P_a(s)(t) $ نمایش داده می‌شود)،
\item
$ V $ به هر گزاره‌ی اتمی مجموعه‌ای از جهان‌ها نسبت می‌دهد
\end{itemize}
مجوعه‌ی همه‌ی مدل‌های کریپکی احتمالاتی را $ \mbbM_{PKL} $ می‌نامیم. برای هر زیرمجموعه‌ی $ E $ از $ S_{a,s} $ از آنجا که مجموعه‌های اندازه‌پذیر زیرمجموعه‌های $ S_{a,s} $ هستند، داریم: 
$$P_a(s)(E)=\sum_{t\in E}P_a(s)(t)$$
و برای هر فرمول $ \varphi $ در زبان تعریف می‌کنیم:
$$P_a(s)(\varphi)=\sum_{\{v\in S_{a,s}\mid M,v\vDash \varphi\}}P_a(s)(t)$$\\
\end{definition}
\begin{definition}\textbf{زبان اعلان عمومی احتمالاتی $ \mL_{PPAL} $.}\index{زبان!اعلان عمومی!احتمالاتی}
این زبان بر پایه‌ی مجموعه‌ی شمارای $ \mbbP $ از گزاره‌های اتمی، مجموعه‌ی متناهی $ \mA $ از عامل‌ها، عملگر کریپکی $ \square_a $، عملگر به‌روزرسانی $ [.] $ و نماد تابعی احتمالاتی $ \mbP_a $ شکل می‌گیرد. فرمول‌های خوش‌تعریف با استفاده از فرم \lr{Backus-Naur} به‌صورت زیر بیان می‌شوند:

\begin{equation*}
\varphi,\psi ::=\ \top\mid\bot\mid p\mid \neg\varphi\mid\varphi\land\psi\mid \square_a\varphi\mid [!\varphi]\psi\mid\sum_{i=1}^n r_i \mbP_a(\varphi_i)\geq r
\end{equation*}
که در آن $ p\in \mbbP $، $ a\in\mA $ و $ r_1,\ldots,r_n,r\in\mbbQ $. $\sum_{i=1}^n r_i \mbP_a(\varphi_i)$  یک نام از زبان $\mL_{PPAL}$ خوانده می‌شود، و $ \sum_{i=1}^n r_i \mbP_a(\varphi_i)\geq r $ فرمول $ -a $احتمالاتی از $\mL_{PPAL}$ نامیده می‌شود. قرار دهید $ \Gamma_{PPAL} $ مجموعه‌ی همه‌ی فرمول‌های خوش‌تعریف و $ P_{a,\mL_{PPAL}} $ مجموعه‌ی همه‌ی فرمول‌های $ -a $احتمالاتی از زبان $\mL_{PPAL}$ باشد، و $ T_{\mL_{PPAL}} $ را نیز مجموعه‌ی همه‌ی نام‌ها قرار دهید. اگر از این زبان عملگر به‌روزرسانی را حذف کنیم آن را $ \mL_{PKL} $ می‌نامیم.

علاوه بر همه‌ی خلاصه‌نویسی‌های پیش‌گفته‌ی قابل بیان در این زبان، از آنجا که عملی غیر از اعلان عمومی در این زبان مطرح نیست برای سادگی نوشتار از $ [\varphi]\psi $ به جای $ [!\varphi]\psi $ به‌عنوان خلاصه‌نویسی استفاده می‌کنیم.
\end{definition}
به منظور تعبیر این زبان می‌بایست به طور همزمان دو تعریف مطرح شود، یکی تعریف راستی و دیگری تعریف مدل‌های به‌روز شده. این دو تعریف به یکدیگر وابسته‌اند ولی به دور نمی‌انجامد.

\begin{definition}\label{semanticK}{\textbf{معناشناسی منطق اعلان عمومی احتمالاتی.}} \index{معناشناسی!اعلان عمومی احتمالاتی}
درستی فرمول $ \varphi\in\Gamma_{\mL_{PPAL}} $ در $ s\in S $ در مدل کریپکی احتمالاتی $ M $ با نماد $ M,s\vDash\varphi $ به‌صورت زیر تعریف می‌شود:
\\

\semanticsa{$  M,s\vDash p $}{$ s\in V(p) $}

\semanticsa{$  M,s\vDash\neg\varphi $}{$ M,s\nvDash\varphi $}

\semanticsa{$  M,s\vDash\varphi\wedge\psi $}{$  M,s\vDash\varphi $ و $  M,s\vDash\psi $}

\semanticsa{$  M,s\vDash \square_a\varphi $}{برای هر $ v\in S $, اگر $ s\xra{a}v $، آنگاه $  M,v\vDash\varphi $}

\semanticsa{$ M,s\vDash[A]\psi $}{$ M|A,s\vDash\psi $ (تعریف \ref{updatedmodel} را ببینید)}

\semanticsa{$  M,s\vDash \sum_{i=1}^n r_i \mbP_a(\varphi_i)\geq r $}{$ \sum_{i=1}^n r_i P_{a,s}(\varphi_i)\geq r $.}\\
\end{definition}
\section{معمای \texorpdfstring{ \lr{Monty Hall}}{Monty Hall}}
\begin{itemize}
\item[]\textiranic{
فرض کنید در یک مسابقه‌ی تلویزیونی شرکت کرده‌اید، و باید از میان سه در یکی را انتخاب کنید با این وصف که پشت یکی از آنها اتومبیل است و پشت دو در دیگر دوچرخه. شما دری را انتخاب می‌کنید، مثلاً در شماره 1، و مجری، که می‌داند پشت هر در چه چیزی نهفته است، دری دیگر را باز می‌کند که پشت آن دوچرخه است، مثلاً در شماره 3. او از شما می‌پرسد «آیا حاضرید دری که انتخاب کرده‌اید را با در شماره 2 عوض کنید؟». سئوال اینجاست که عوض کردن در به نفع شماست یا نه؟
}
\end{itemize}

همانطور که در \citep{Kooi2003} آمده است خانم سَوِنت\LTRfootnote{Savant} که مرتباً در کتاب گینس به‌عنوان باهوش‌ترین فرد رکورد داشته است بر این باور است که اگر انتخاب را تغییر دهید در یک سوم موارد دوچرخه می‌برید و در دو سوم موارد اتومبیل. او اینگونه استدلال می‌کند که فرض کنید شما در مرحله‌ی اول دری را انتخاب کردید که پشت آن اتومبیل است، بنابر این شما نباید انتخابتان را تغییر دهید و این در یک سوم موارد اتفاق می‌افتد. از طرف دیگر فرض کنید که انتخاب اولیه‌ی شما دری باشد که پشتش دوچرخه است، که در دو سوم موارد رخ می‌دهد. مجری نمی‌تواند دری که پشتش اتومبیل است و دری که شما انتخاب کرده‌اید را باز کند او مجبور است در دیگری را که پشتش دوچرخه است باز کند. بنابراین در حالتی که شما ابتدا دری را انتخاب کرده‌اید که پشتش دوچرخه است تغییر انتخاب، بردن ماشین را تضمین می‌کند. پس با تغییر انتخاب در دو سوم موارد برنده‌ی اتومبیل خواهید شد.

می‌خواهیم به کمک اثباتی صوری در \lr{PPAL} نشان ‌دهیم که تعویض در به نفع شماست و مدعای خانم سَوِنت را اثبات کنیم.

قبل از آن به چند لم نیازمندیم که آنها را در اینجا اثبات می‌کنیم.

\begin{lemma}\label{negP}
$ \mbP_a(\varphi)=1-\mbP_a(\neg\varphi) $
\end{lemma}
\bp
از اصل جمع‌پذیری متناهی داریم
$ \mbP_a(\top\wedge\varphi)+\mbP_a(\top\wedge\neg\varphi)=\mbP_a(\top) $
و با استفاده از احتمال راستی حکم برقرار است.
\ep
\begin{lemma}\label{identity2}
در فرمول $ \sum_{i=1}^n r_i\mbP_a(\varphi_i)\geq r $، اگر داشته باشیم   $  \mbP_a(\varphi_j)=r' $، که در آن $ 1\leq j\leq n $ و $ r' $ عددی گویا باشد، آنگاه می‌توان اثبات کرد:
$$\sum_{i=1 , i\neq j}^n r_i\mbP_a(\varphi_i)\geq r-r'$$
\end{lemma}
\bp
از فرضیات و با استفاده از خلاصه‌نویسی فرمول‌ها و اصل $\0 $-نام‌ها داریم
$$r_1\mbP_a(\varphi_1) +\cdots +r_j\mbP_a(\varphi_j) +\cdots+r_n\mbP_a(\varphi_n)\geq r$$
$$\0\mbP_a(\varphi_1)+\cdots-\mbP_a(\varphi_j) +\cdots+\0\mbP_a(\varphi_n)=-r'$$
سپس با استفاده از اصل افزودن، حکم قضیه اثبات می‌شود.
\ep
\begin{lemma}\label{identity3}
در فرمول $ [\varphi]\sum_{i=1}^n r_i\mbP_a(\varphi_i)\geq r $، اگر داشته باشیم   $  [\varphi]\mbP_a(\varphi_j)=r' $، که در آن $ 1\leq j\leq n $ و $ r' $ عددی گویا باشد، آنگاه می‌توان اثبات کرد:
$$[\varphi]\sum_{i=1}^n r_i\mbP_a(\varphi_i)\geq r\leftrightarrow[\varphi]\sum_{i=1,i\neq j}^n r_i\mbP_a(\varphi_i)\geq r-r_j r'$$
\end{lemma}
\bp
فرض کنید $ \mbP_a(\varphi)>\0 $ آنگاه هم‌ارزی‌های زیر با استفاده از اصل به‌روزرسانی احتمال 1 برقرارند:
$$[\varphi]\sum_{i=1}^n r_i\mbP_a(\varphi_i)\geq r\leftrightarrow\sum_{i=1}^n r_i\mbP_a(\varphi\wedge[\varphi]\varphi_i)\geq r\mbP_a(\varphi)$$
$$[\varphi]\mbP_a(\varphi_j)=r'\leftrightarrow\mbP_a(\varphi\wedge[\varphi]\varphi_j)=r'\mbP_a(\varphi)$$

حال با استفاده از لم \ref{identity} و خلاصه‌نویسی فرمول‌ها بدست می‌آوریم
$$[\varphi]\sum_{i=1}^n r_i\mbP_a(\varphi_i)\geq r\leftrightarrow\sum_{i=1,i\neq j}^n r_i\mbP_a(\varphi\wedge[\varphi]\varphi_i)\geq (r-r_j r')\mbP_a(\varphi)$$

و در نتیجه اصل به‌روزرسانی احتمال 1 حکم را نتیجه می‌دهد. برای حالت $ \mbP_a(\varphi)=\0 $ نیز مشابه همین استدلال با کمک اصل به‌روزرسانی احتمال 2 برقرار است.
\ep

اکنون به مدل‌سازی معما در دستگاه منطقی‌مان می‌پردازیم و سپس به اثبات مدعا روی می‌آوریم.

مجموعه‌ی عامل‌ها را $ \mA=\{c,m\} $ می‌گیریم ($ c $ معرف شرکت‌کننده و $ m $ معرف مجری) و مجموعه‌ی گزاره‌های اتمی را اجتماع سه مجموعه‌ی $ A $، $ C $ و $ O $. بطوریکه\\
$ A=\{A_1,A_2,A_3\} $ که در آن $ A_i $ یعنی اتومبیل پشت در شماره $ i $ است.\\
$ C=\{C_1,C_2,C_3\} $ که در آن $ C_i $ یعنی شرکت‌کننده ابتدا در شماره $ i $ را انتخاب کرده است.\\
$ O=\{O_1,O_2,O_3\} $ که در آن $ O_i $ یعنی در شماره $ i $ توسط مجری باز شده است.

اکنون قواعد بازی را مدل می‌کنیم. این قواعد به این شرح هستند که: تنها یک اتومبیل پشت درهاست، شرکت‌کننده تنها می‌تواند یک در را انتخاب کند و مجری تنها می‌تواند یک در را باز کند.
$$onecar\ =\oplus A,\quad onechoice\ =\oplus C,\quad oneopen\ =\oplus O$$
که $ \oplus $ یعنی «یای انحصاری\LTRfootnote{\lr{exclusive or}}». فرض می‌کنیم که شرکت‌کننده می‌بایست به اینکه اتومبیل پشت دری خاص قرار دارد احتمال $ \frac{1}{3} $ نسبت دهد. همجنین فرض می‌کنیم شرکت‌کننده با انتخاب یک در چیزی در مورد جایگاه اتومبیل کشف نمی‌کند. بنابراین شرکت‌کننده بعد از انتخاب در نیز می‌بایست این احتمال را همان $ \frac{1}{3} $ در نظر بگیرد، یعنی انتخاب شرکت‌کننده مستقل است از جایی که اتومبیل قرار دارد.
$$equal\ =\!\!\!\!\!\bigwedge_{i\in \{1,2,3\}}\!\!\!\!\!\mbP_c(A_i)=\frac{\1}{\3},\qquad\qquad independentAC\ =\!\!\!\!\!\bigwedge_{j\in\{1,2,3\}}\!\!\!\!\![C_j]equal$$

بخش اساسی بررسی این معما آن است که ببینیم تحت چه شرایطی مجری دری را باز می‌کند. او دقیقاً یک در را باز می‌کند به شرطی که شرکت‌کننده آن را انتخاب نکرده باشد و اتومبیل نیز پشت آن نباشد.
$$conditions=\!\!\!\!\!\bigwedge_{i,j\in \{1,2,3\}}[C_i](O_j\leftrightarrow(\neg A_j\wedge\neg C_j\wedge\!\!\!\bigwedge_{k\in\{1,2,3\},k\neq j}\!\!\!\!\!\!\!\!\neg O_k))$$

حال قرار دهید
$$initial=\ onecar\wedge onechoice\wedge oneopen\wedge equal\wedge independentAC\wedge conditions$$

سئوال این است که شرکت‌کننده انتخاب خود را تغییر دهد یا نه:
$$switch=\ [C_1][O_3]\mbP_c(A_1)\leq \mbP_c(A_2)$$

اگر این جمله درست باشد، احتمال اینکه شرکت‌کننده ماشین را ببرد با تغییر در انتخاب کاهش نمی‌یابد. معلوم می‌شود که $ initial $ برای بدست آوردن این نتیجه کفایت نمی‌کند. آنچه ضروری است آن است که شرکت‌کننده از برقراری شرایط اولیه‌ی بازی مطمئن باشد :$ \mbP_c(initial)=\1 $. ما همچنین به دو فرض طبیعی نیز نیاز داریم، اولاً از نظر شرکت‌کننده احتمال اینکه او در شماره 1 را انتخاب کند بزرگتر از صفر است: $ \mbP_c(C_1)>\0 $ ثانیاً بعد از اینکه شرکت‌کننده در شماره 1 را انتخاب کرد از نظر او احتمال اینکه مجری در 3 را باز کند بزرگتر از صفر است: $ [C_1]\mbP_c(O_3)>\0 $. اینها برای بدست آوردن $switch$ کفایت می‌کنند.

فرض $independentAc$ دلالت دارد بر اینکه $ [C_1]\mbP_c(A_1)=\frac{\1}{\3} $ و بنابراین:
\begin{equation}\label{a}
\mbP_c(A_1)=\mbP_c(O_3\wedge A_1)+\mbP_c(\neg O_3\wedge A_1)\ \Rightarrow\ [C_1]\mbP_c(O_3\wedge A_1)\leq\frac{\1}{\3}
\end{equation}
با در نظر گرفتن شرایط $conditions$،  $onechoice$ و $onecar$ خواهیم داشت:
\begin{equation}\label{b0}
 [C_1]\mbP_c(A_2\rightarrow O_3)=\1
\end{equation}
زیرا از شرایط $onechoice$ و $onecar$ داریم
$$C_1\leftrightarrow\neg C_1\wedge\neg C_3\quad\textrm{و}\quad A_2\leftrightarrow\neg A_1\wedge\neg A_3$$
و در نتیجه 
\begin{equation}\tag{1}
C_1\wedge A_2\rightarrow \neg C_3\wedge\neg A_3
\end{equation}
از طرف دیگر بنابر شرط $ conditions $ داریم $ [C_1]O_1\rightarrow[C_1]\neg C_1 $. از این و از $ C_1\rightarrow[C_1]C_1 $ نتیجه می‌شود $ C_1\wedge O_1\rightarrow\bot $ که معادل است با
\begin{equation}\tag{2}
\neg C_1\vee \neg O_1
\end{equation}
و به همین صورت داریم
\begin{equation}\tag{3}
\neg A_2\vee \neg O_2
\end{equation}
پس با استفاده از (1)، (2) و (3) بدست می‌آید
$$C_1\wedge A_2\rightarrow \neg C_3\wedge\neg A_3\wedge\neg O_1\wedge\neg O_2$$
و از قاعده‌ی ضرورت $ [C_1] $ داریم
$$[C_1](C_1\wedge A_2)\rightarrow [C_1](\neg C_3\wedge\neg A_3\wedge\neg O_1\wedge\neg O_2)$$
و از شرط $ conditions $ نتیجه می‌شود $ [C_1](C_1\wedge A_2)\rightarrow [C_1]O_3 $ که به خاطر اصل ثبات اتم معادل است با
$$C_1\rightarrow(A_2\rightarrow O_3)$$
در نتیجه داریم $ [C_1]C_1\rightarrow(A_2\rightarrow O_3) $ و بنابراین $ C_1\rightarrow[C_1](A_2\rightarrow O_3) $ و
$$C_1\leftrightarrow C_1\wedge[C_1](A_2\rightarrow O_3)$$
و با استفاده از قاعده‌ی هم‌ارزی نتیجه می‌شود
$$\mbP_c(C_1\wedge[C_1](A_2\rightarrow O_3))=\mbP_c(C_1)$$
و چون طبق اصل به‌روزرسانی-احتمال 1 داریم
$$\mbP_c(C_1)>\0\rightarrow([C_1]\mbP_c(A_2\rightarrow O_3)=\1\leftrightarrow \mbP_c(C_1\wedge[C_1](A_2\rightarrow O_3))=\mbP_c(C_1))$$
گزاره‌ای که به دنبالش بودیم اثبات می‌شود.
\\ \\
حال با استفاده از \ref{b0} بدست می‌آید:
\begin{equation}\label{b}
 [C_1]\mbP_c(O_3\wedge A_2)=\mbP_c(A_2) 
\end{equation}
زیرا از لم \ref{negP} می‌دانیم
$$ \mbP_c(A_2\wedge\neg O_3)=\0 $$
و از اصل جمع‌پذیری متناهی داریم
$$ \mbP_c(O_3\wedge A_2)+\mbP_c(\neg O_3\wedge A_2)=\mbP_c(A_2) $$
پس با استفاده از لم \ref{identity2} بدست می‌آید
$$ \mbP_c(O_3\wedge A_2)=\mbP_c(A_2) $$
و سپس قاعده‌ی ضرورت $ [C_1] $ حکم را اثبات می‌کند.

(\ref{b})
 به همراه $ [C_1]\mbP_c(A_2)=\frac{\1}{\3} $
(که از $independentAC$ بدست می‌آید) و با استفاده از لم \ref{identity3} ما را مجاز می‌کند که نتیجه بگیریم $ [C_1]\mbP_c(O_3\wedge A_2)=\frac{\1}{\3} $. حال خلاصه‌نویسی فرمول‌ها و لم \ref{identity3} (قرار دهید $ r_j=1 $، $ r'=\frac{\1}{\3} $ و $ r=0 $) و (\ref{a}) نتیجه دهد:
$$[C_1]\mbP_c(O_3\wedge A_1)\leq\mbP_c(O_3\wedge A_2)$$
و با استفاده از اصل ثبات اتم خواهیم داشت
\begin{equation}\label{c}
[C_1]\mbP_c(O_3\wedge[O_3]A_1)\leq\mbP_c(O_3\wedge[O_3]A_2)
\end{equation}
و در نهایت اثبات می‌شود
\begin{equation}\tag{switch}
[C_1][O_3]\mbP_c(A_1)\leq\mbP_c(A_2)
\end{equation}
\clearpage{\pagestyle{empty}\cleardoublepage}
\backmatter
{\small
{\baselineskip=.75cm
\addcontentsline{toc}{chapter}{واژه‌نامه فارسی به انگلیسی}
\thispagestyle{empty}
\chapter*{واژه‌نامه فارسی به انگلیسی}
\markboth{واژه‌نامه فارسی به انگلیسی}{واژه‌نامه فارسی به انگلیسی}

\noindent
\englishgloss{probability}{احتمال}
\englishgloss{posterior probability}{احتمال پسینی}
\englishgloss{prior probability}{احتمال پیشینی}
\englishgloss{occurrence probability}{احتمال رخداد}
\englishgloss{propositional probability}{احتمال گزاره‌ای}
\englishgloss{observation probability}{احتمال {\observation}}
\englishgloss{axiom}{اصل موضوع}
\englishgloss{reduction axioms}{اصول موضوعه‌ی {\reduction}}
\englishgloss{public announcement}{اعلان عمومی}
\englishgloss{probability measure}{اندازه‌ی احتمالاتی}
\englishgloss{inner probability measure}{اندازه‌ی احتمالاتی درونی}
\englishgloss{static}{ایستا}
\englishgloss{belief}{باور}
\englishgloss{update}{به‌روزرسانی}
\clearpage{\pagestyle{empty}\cleardoublepage}
\addcontentsline{toc}{chapter}{واژه‌نامه  انگلیسی به  فارسی}
\thispagestyle{empty}
\chapter*{واژه‌نامه  انگلیسی به  فارسی}
\markboth{واژه‌نامه  انگلیسی به  فارسی}{واژه‌نامه  انگلیسی به  فارسی}

\noindent
\persiangloss{عامل}{agent}
\persiangloss{اصل موضوع}{axiom}
\persiangloss{باور}{belief}
\persiangloss{همه‌دانی مشترک}{common knowledge}
\persiangloss{تمامیت}{completeness}
\persiangloss{سازگار}{consistent}
\persiangloss{جمع‌پذیر شمارا}{countably additive}
\persiangloss{تمییز دادن}{distinguish}
\persiangloss{پویا}{dynamic}
\persiangloss{شناخت}{epistemic}
\persiangloss{عمل}{event}
\persiangloss{یای انحصاری}{exclusive or}
\persiangloss{جمع‌پذیر متناهی}{finitely additive}
\persiangloss{فرمول}{formula}
\persiangloss{فرادانش}{higher order information}
\persiangloss{خودبیمار انگار}{hypochondriac}
\persiangloss{اندازه‌ی احتمالاتی درونی}{inner probability measure}
\persiangloss{خودآگاهی}{intropection}}
\clearpage{\pagestyle{empty}\cleardoublepage}
{\baselineskip=.6cm
\phantomsection
\addcontentsline{toc}{chapter}{نمایه}
\printindex}
\clearpage{\pagestyle{empty}\cleardoublepage}
\phantomsection
\addcontentsline{toc}{chapter}{مراجع}
\bibliographystyle{chicago-fa}
\bibliography{MyReferences}
\clearpage{\pagestyle{empty}\cleardoublepage}
}
%ایجاد صفحه‌ای سفید%%%%%%%%%%%%%%%%%%%%%%%%%%%%%%%%%%%
\newpage
\thispagestyle{empty}
\clearpage
~~~
%%%%%%%%%%%%%%%%%%%%%%%%%%%%%%%%%%%%
% در این فایل، عنوان پایان‌نامه، مشخصات خود و چکیده پایان‌نامه را به انگلیسی، وارد کنید.
% توجه داشته باشید که جدول حاوی مشخصات پایان‌نامه/رساله، به طور خودکار، رسم می‌شود.
%%%%%%%%%%%%%%%%%%%%%%%%%%%%%%%%%%%%
\baselineskip=.6cm
\begin{latin}
\latinuniversity{Shahid Beheshti University}
\latinfaculty{Faculty of Mathematics & Computer Science}
\latindegree{M. Sc. }
%group:
\latinsubject{Department of Computer Science}
\latinfield{Computer Sciences}
\latintitle{Concurency Control in Native XML Databases}
\firstlatinsupervisor{Dr. Sayyed Kamyar Izadi}
%\secondlatinsupervisor{Second Supervisor}
\firstlatinadvisor{Dr. Mahmoud Fazlali}
%\secondlatinadvisor{Second Advisor}
\latinname{Mohsen}
\latinsurname{Lachinani}
\latinthesisdate{2015}
\latinkeywords{Native XML , Concurrency Control , Query Processing}
\en-abstract{\noindent
In this thesis, we first introduce an instance of probabilistic epistemic logics (PEL) and prove its completeness. Then in our approach toward generalization we will consider a simple model of this logic to develop it to dynamic logics which are able to model information changes in multi-agent systems.
\\
After a short description of non-probabilistic dynamic epistemic logics, we also introduce probabilistic dynamic epistemic logic (PDEL) by taking into account three sources of probability, namely, prior probability of states, occurrence probability for events based on a process corresponding to agents’ view, and the probability of uncertainty of observing events. This three sources are used to provide a generalized update mechanism that is a natural and convenient format for modeling information flow.
\\
Then in order to prove a completeness of probabilistic dynamic epistemic logic from the completeness of probabilistic epistemic logic we present axioms which are sound to reduce the formulas containing dynamic operators to the formulas in the corresponding static language.
\\
Finally we will introduce a kind of probabilistic dynamic epistemic logic which is adapted to solve the Monty Hall dilemma and will prove its completeness. Then we will obtain a formal solution for this dilemma within this logic.
}
%\latinvtitle
\end{latin}
\label{LastPage}
\end{document}